\part{Foundations}

\chapter{Why build your own model}
\chapfig{fig_spark.png}

For most companies, ``using AI'' means calling someone else's API. That is the right default. So why would
an enterprise ever train a language model \emph{from scratch}?

There are four honest reasons, and they are the reasons this project exists.

\section{Control, privacy, and provenance}
A model you trained is a model you own. Every weight has a known origin; no data left your control; nothing
about its behaviour is a black box you rent by the token. In regulated domains --- healthcare, finance, law
--- ``we can show exactly what this model was trained on'' is not a nice-to-have, it is the difference
between deployable and not.

\section{A tokenizer that speaks your language}
This is the subtle one, and it is the single capability you \emph{cannot} get by fine-tuning someone else's
model. General tokenizers shatter domain words: \texttt{pharmacokinetics} or \texttt{indemnification} become
five or six tokens, wasting the model's budget on spelling rather than meaning. When you train from scratch
you build the vocabulary on day one. Ours reached a \textbf{fertility of 1.57} tokens per word on dense
pharmaceutical text --- roughly half what an off-the-shelf tokenizer would use (Chapter~4).

\section{Cost at inference scale}
A small specialist that runs on a cheap GPU can be dramatically cheaper to serve than a frontier API once
you are doing millions of calls --- and it never changes underneath you.

\section{Understanding}
You cannot truly debug, trust, or improve a system you did not build. The team that trains a model from
scratch understands tokenization, attention, learning-rate schedules, and evaluation in a way no amount of
prompt engineering will teach.

\begin{numbers}[What ``from scratch'' bought us, concretely]
A 350M model that writes fluent pharmaceutical prose; a custom 32k tokenizer; a fully transparent training
recipe; and --- when we needed it --- the ability to swap the entire architecture and scale up, because we
owned every line. Total compute for the whole journey: roughly \textbf{\$1{,}500} across twelve models.
\end{numbers}

\section{When \emph{not} to}
Be honest: if your task is well served by an API, or you lack a large, clean domain corpus, or you need
frontier-level reasoning, building from scratch is the wrong call. From scratch wins when you have abundant
domain text, when the tokenizer matters, and when control and cost-at-scale dominate. Pharmaceuticals --- with
PubMed's tens of billions of open words --- is a near-perfect fit, which is exactly why the best small
biomedical models (BioGPT, PubMedBERT, BioMedLM) were all trained this way.

\begin{lesson}
Build from scratch when you have (a) a large clean domain corpus, (b) a real need for a custom tokenizer or
full provenance, and (c) a task that plays to a small model's strengths. Otherwise, fine-tune or call an API.
\end{lesson}

\chapter{The recipe, end to end}
\chapfig{fig_pipeline.png}

A language model is six honest steps. The whole rest of this book is these six steps done carefully, plus
the engineering and economics around them.

\begin{enumerate}[leftmargin=2.2em]
  \item \textbf{Gather data.} Blend a large general-English corpus with your domain text. Decide the mix.
  \item \textbf{Train a tokenizer.} A custom byte-level BPE over the blended corpus.
  \item \textbf{Define the model.} A modern decoder-only Transformer, written by hand.
  \item \textbf{Pretrain.} Predict the next token, billions of times, across many GPUs.
  \item \textbf{Instruction-tune (SFT).} Teach the raw next-token engine to \emph{answer}.
  \item \textbf{Evaluate.} Grade it on real held-out tasks, and read its answers.
\end{enumerate}

\fig{figures/diag_pipeline6.png}{0.98}{The six stages, in order. Everything in this book is one of these
boxes done carefully --- plus the engineering and economics that surround them.}

\noindent Three ideas thread through all six and separate a toy from a real model:

\subsection{Learn the domain \emph{and} the language}
The famous medical models (Meditron, SaulLM) use 99\% domain text --- but those are \emph{continued}
pretraining on a base that already speaks English. From scratch, you must teach English too, so we mix far
more general text (Chapter~3). Getting this ratio right is the project's central data decision.

\subsection{Over-train the small model}
Chinchilla's ``20 tokens per parameter'' is a compute-optimal \emph{floor}, not a target. Every good small
model today is trained far past it, because that is what turns ``grammatical'' into ``fluent.'' We trained a
350M model on 20 billion tokens ($\approx$57$\times$).

\subsection{Judge on tasks, not on loss}
The most important lesson in the book, learned painfully in Chapter~7: the model with the \emph{lowest training
loss} gave the \emph{worst answers}. Lower perplexity does not mean better. Always measure on the real task.

\begin{lesson}
The pipeline is simple; the judgement is not. The three decisions that matter most are the data mix, how long
to train, and \emph{what you measure}. Get those right and the rest is engineering.
\end{lesson}
